\subsection{Experimental demonstrations regarding TDI and clock noise transfer}
\subsubsection{JPL Interferometry Testbed (de Vine et al., 2010)}

The first reported hardware demonstration of TDI was carried out at JPL by de Vine et al. [REF]. The setup implemented a two-spacecraft configuration with two lasers per spacecraft and independent phasemeter clocks and static arm lengths of about 1 m — reproducing several of the key experimental complexities of LISA including multiple heterodyne frequencies, independent clocks, and independent phase measurements of both optical signals and clock sidebands. White noise was injected into the inter-spacecraft phase-locked loop to bring laser frequency noise to approximately 800 Hz/$\sqrt{\text{Hz}}$. The experiment demonstrated TDI suppression of laser frequency noise by roughly nine orders of magnitude and clock phase noise by about four to five orders of magnitude at 3 mHz, recovering the intrinsic displacement noise floor of the testbed at around $10^{-5}$ cycles/$\sqrt{\text{Hz}}$. One limitation was that the arm lengths were static and short, and physical light-travel times representative of the LISA constellation were not reproduced.

\subsubsection{UFLIS (Mitryk et al., 2010 and 2012)}
The University of Florida Laser Interferometry Simulator (UFLIS) was the first setup to include electronic phase delay units capable of time-varying delays on the order of seconds, allowing spacecraft motion effects to be included for the first time in a hardware TDI testbed [REF, REF]. Each spacecraft was represented by one laser: two were PDH-locked to ULE cavities, achieving a combined beatnote stability of 30–300 Hz/$\sqrt{\text{Hz}}$ near LISA pre-stabilization levels, while the remaining lasers were phase-locked to this cavity-stabilized reference with MHz offsets. The delays were applied by the Electronic Phase Delay (EPD) unit with delay rates updated continuously during a run to simulate time-changing inter-spacecraft separations. UFLIS showed suppression of pre-stabilized laser phase noise to approximately $10^{-5}$ cycles/$\sqrt{\text{Hz}}$, recovering the testbed's intrinsic noise floor, and demonstrated that second-generation TDI outperforms first-generation combinations when delays are linearly changing. GW signal injection was also demonstrated, with sinusoidal signals recovered in the TDI spectrum. Its primary limitation relative to the present work is the absence of independent clock sources: all spacecraft clocks were synchronized, meaning clock noise and its suppression via sideband transfer could not be studied.

\subsubsection{Hexagon / AEI Clock Synchronization (Yamamoto et al., 2022)}
The Hexagon interferometer at the Albert Einstein Institute in Hannover is a picometer-stable optical testbed for the LISA phasemeter, extended to simulate inter-spacecraft clock transfer by equipping each simulated spacecraft with an independent clock and GHz sideband modulation via EOMs [REF, REF]. Without clock synchronization, the three-signal combination is dominated by differential clock jitter. After applying clock synchronization via TDIR-like processing and sideband measurements, the clock noise was suppressed by three orders of magnitude at 1 Hz and up to six orders of magnitude at 0.1 mHz, bringing the residual below the LISA single-link requirement of 10 pm/$\sqrt{\text{Hz}}$ across the observation band. The post-correction floor is set by the optical bench sensitivity, which is close to but not yet at the LISA 1 pm/$\sqrt{\text{Hz}}$ single-link mark, with ongoing noise hunting to close the remaining gap. The experiment also provided the first hardware observation of several secondary noise couplings relevant to TDI: \ac{ADC} aliasing, the flexing-filtering coupling, and Lagrange interpolation error. The Hexagon does not implement programmable inter-spacecraft delays, and is therefore complementary to testbeds focused on TDI laser noise suppression. Work on Hexagon continues as it is a part of the official LISA test plan.

\subsubsection{LISA on Table / LOT (Vidal 2023, 2025)}
The LISA On Table (LOT) experiment, developed at the AstroParticule et Cosmologie (APC) laboratory in Paris, is an electro-optical simulator in which laser frequency noise and inter-spacecraft delays are generated entirely by \ac{DDS}, with no physical optical path [REF]. Clock noise and GW signal injection are not included. With DDS-generated laser noise of approximately 100$\sqrt{2}$ Hz/$\sqrt{\text{Hz}}$, both TDI 1.0 and 2.0 suppress laser frequency noise to the intrinsic testbed floor of $10^{-5}$ cycles/$\sqrt{\text{Hz}}$, consistent with other testbeds. Beyond laser noise suppression, the LOT was used to experimentally identify and characterize the non-commutativity between the delay operator and the phasemeter CIC decimation filters, producing residual aliased laser noise with a characteristic f² slope — a finding validated against an analytical model and shown to require at least a third-order CIC filter in LISA to remain below the allocated phasemeter noise level [REF].

\subsubsection{TianQin Clock Transfer (Zeng et al., 2023)}
The TianQin group at Sun Yat-sen University demonstrated inter-spacecraft clock jitter readout based on laser sideband phase transfer under weak-light conditions ($\approx$1~nW received power), with a simulated $\approx$50~µs propagation delay using a 10~km optical fiber [REF]. Each simulated spacecraft carried an independent clock — one oven-controlled crystal oscillator and one rubidium clock — with GHz sideband modulation via EOMs following the same clock transfer scheme as LISA. Comparison between electrical and optical clock transfer paths identified the GHz frequency synthesizer as the dominant noise source below 50~mHz. After differential clock noise correction, the residual optical chain noise was below 40~fs/$\sqrt{\text{Hz}}$ above 6~mHz, and carrier phase accuracy reached $10^{-4}$~cycles/$\sqrt{\text{Hz}}$ above 6~mHz, recovering the phasemeter noise floor. The experiment does not include programmable arm-length delays, TDI post-processing, or laser frequency noise, and thus addresses only the clock transfer component of the full metrology chain.

\subsubsection{HUST Picometer Signal Recovery (Xu et al., 2023)}
Xu et al.\ demonstrated combined TDI-like laser noise cancellation and clock noise removal with physical signal recovery [REF]. The setup used a single laser split through two AOMs at 75~MHz and 85~MHz to create uncorrelated effective laser noises between two simulated spacecraft, each carrying an independent clock with GHz sideband modulation via EOMs. A PZT-actuated mirror injected a 60~pm, 1~Hz displacement signal. With input laser frequency noise at approximately $10^{3}$~cycles/$\sqrt{\text{Hz}}$ and clock noise several orders of magnitude above the noise floor, the post-processing suppressed laser frequency noise by five orders of magnitude and clock noise by two orders of magnitude, recovering the 60~pm signal cleanly at the intrinsic interferometer floor of $\sim$$10^{-11}$~m/$\sqrt{\text{Hz}}$ at 1~Hz. Notably, no arm-length delay is implemented — the equal-arm geometry cancels laser noise directly — so the experiment validates signal recovery under the TDI equal-arm principle rather than testing TDI under realistic unequal delays. It does not address time-varying inter-spacecraft separations, ranging, or LISA-band performance below 1~mHz.